% 数学建模竞赛论文 LaTeX 导言区
% 使用方法：在 R Markdown YAML 中添加 includes: in_header: preamble.tex

% ============================================
% 页面设置
% ============================================
\usepackage{geometry}
\geometry{
  a4paper,
  top=2.5cm,
  bottom=2.5cm,
  left=2.5cm,
  right=2.5cm
}

% ============================================
% 页眉页脚
% ============================================
\usepackage{fancyhdr}
\pagestyle{fancy}
\fancyhf{}
\fancyhead[C]{\small 数学建模竞赛论文}
\fancyfoot[C]{\thepage}
\renewcommand{\headrulewidth}{0.4pt}

% ============================================
% 数学环境
% ============================================
\usepackage{amsmath, amssymb, amsthm}
\usepackage{bm}

% 定理环境
\newtheorem{theorem}{定理}[section]
\newtheorem{lemma}[theorem]{引理}
\newtheorem{corollary}[theorem]{推论}
\newtheorem{proposition}[theorem]{命题}

% 定义环境
\newtheorem{definition}[theorem]{定义}
\newtheorem{example}[theorem]{例}

% 假设环境
\newtheorem{assumption}{假设}

% ============================================
% 表格
% ============================================
\usepackage{booktabs}
\usepackage{multirow}
\usepackage{array}

% 表格列类型
\newcolumntype{C}[1]{>{\centering\arraybackslash}p{#1}}
\newcolumntype{L}[1]{>{\raggedright\arraybackslash}p{#1}}
\newcolumntype{R}[1]{>{\raggedleft\arraybackslash}p{#1}}

% ============================================
% 图片
% ============================================
\usepackage{graphicx}
\graphicspath{{../figures/}}

% 图片位置控制
\usepackage{float}
\usepackage[figurewithin=section]{caption}

% ============================================
% 代码高亮
% ============================================
\usepackage{listings}
\usepackage{xcolor}

\lstset{
  language=R,
  basicstyle=\ttfamily\small,
  keywordstyle=\color{blue}\bfseries,
  commentstyle=\color{green!60!black}\itshape,
  stringstyle=\color{red},
  numbers=left,
  numberstyle=\tiny\color{gray},
  numbersep=5pt,
  frame=single,
  framesep=5pt,
  breaklines=true,
  breakatwhitespace=true,
  showstringspaces=false,
  tabsize=2,
  captionpos=b,
  backgroundcolor=\color{gray!10}
}

% ============================================
% 算法伪代码
% ============================================
\usepackage{algorithm}
\usepackage{algpseudocode}

\algnewcommand\algorithmicinput{\textbf{输入：}}
\algnewcommand\Input{\item[\algorithmicinput]}
\algnewcommand\algorithmicoutput{\textbf{输出：}}
\algnewcommand\Output{\item[\algorithmicoutput]}

% ============================================
% 超链接
% ============================================
\usepackage{hyperref}
\hypersetup{
  colorlinks=true,
  linkcolor=blue,
  filecolor=magenta,
  urlcolor=cyan,
  citecolor=blue
}

% ============================================
% 中文支持
% ============================================
\usepackage{ctex}

% 中文字体设置（可选）
% \setCJKmainfont{SimSun}[BoldFont=SimHei, ItalicFont=KaiTi]
% \setCJKsansfont{SimHei}
% \setCJKmonofont{FangSong}

% ============================================
% 其他常用宏包
% ============================================
\usepackage{enumitem}  % 列表定制
\usepackage{indentfirst}  % 首行缩进
\usepackage{setspace}  % 行距控制
\usepackage{soul}  % 文字装饰
\usepackage{ulem}  % 下划线

% ============================================
% 自定义命令
% ============================================
% 数学符号
\newcommand{\R}{\mathbb{R}}  % 实数集
\newcommand{\N}{\mathbb{N}}  % 自然数集
\newcommand{\Z}{\mathbb{Z}}  % 整数集
\newcommand{\Q}{\mathbb{Q}}  % 有理数集
\newcommand{\E}{\mathbb{E}}  % 期望
\newcommand{\Var}{\mathrm{Var}}  % 方差
\newcommand{\Cov}{\mathrm{Cov}}  % 协方差
\newcommand{\Prob}{\mathbb{P}}  % 概率

% 优化问题
\newcommand{\argmin}{\operatornamewithlimits{argmin}}
\newcommand{\argmax}{\operatornamewithlimits{argmax}}
\newcommand{\st}{\text{s.t.}}  % 约束条件

% 常用符号
\newcommand{\dd}{\mathrm{d}}  % 微分符号
\newcommand{\pp}{\partial}  % 偏导符号

% ============================================
% 文档信息
% ============================================
% 标题格式
\usepackage{titlesec}
\titleformat{\section}{\Large\bfseries}{\thesection}{1em}{}
\titleformat{\subsection}{\large\bfseries}{\thesubsection}{1em}{}
\titleformat{\subsubsection}{\normalsize\bfseries}{\thesubsubsection}{1em}{}

% 目录设置
\usepackage[contents={目录}, title={目录}]{tocloft}
